\documentclass[11pt]{article}
\usepackage[letterpaper,margin=0.68in,headheight=15pt]{geometry}
\usepackage[T1]{fontenc}
\usepackage{lmodern}
\usepackage{amsmath,amssymb}
\usepackage{array,booktabs,enumitem,multicol}
\usepackage{xcolor}
\usepackage{tikz}
\usepackage{fancyhdr}
\usepackage{microtype}
\setlength{\parindent}{0pt}
\setlength{\parskip}{5pt}
\definecolor{olive}{HTML}{555B35}
\definecolor{gold}{HTML}{B18A22}
\definecolor{soft}{HTML}{F4F1DE}
\definecolor{bluegray}{HTML}{E8EEF2}
\pagestyle{fancy}
\fancyhf{}
\fancyhead[L]{\textcolor{olive}{\small MATH\& 141: Precalculus I}}
\fancyhead[R]{\textcolor{gold}{\small Fall 2026}}
\fancyfoot[C]{\thepage}
\renewcommand{\headrulewidth}{0.45pt}
\renewcommand{\footrulewidth}{0pt}
\setlist[enumerate]{leftmargin=*,itemsep=0.08in,topsep=0.05in}
\newcommand{\assessmenttitle}[2]{%
  {\Large\bfseries Module #1: #2}\par\vspace{3pt}
}
\newcommand{\studentline}{%
\noindent\textbf{Name:}\hspace{2.6in}\textbf{Date:}\hspace{1.15in}\par\vspace{7pt}}
\newcommand{\directions}[1]{\textit{#1}\par\vspace{4pt}}
\newcommand{\answerblank}[1]{\vspace{#1}}
\newcommand{\blankgrid}[4]{%
\begin{center}
\begin{tikzpicture}[x=0.75cm,y=0.75cm]
  \draw[step=1,gray!28,very thin] (#1,#3) grid (#2,#4);
  \draw[->,thick] (#1,0)--(#2,0) node[right] {$x$};
  \draw[->,thick] (0,#3)--(0,#4) node[above] {$y$};
  \foreach \x in {#1,...,#2}{\ifnum\x=0\else\draw (\x,0.10)--(\x,-0.10) node[below=2pt,font=\scriptsize]{\x};\fi}
  \foreach \y in {#3,...,#4}{\ifnum\y=0\else\draw (0.10,\y)--(-0.10,\y) node[left=2pt,font=\scriptsize]{\y};\fi}
\end{tikzpicture}
\end{center}}
\newcommand{\smallgrid}[4]{%
\begin{tikzpicture}[x=0.50cm,y=0.50cm]
  \draw[step=1,gray!28,very thin] (#1,#3) grid (#2,#4);
  \draw[->,thick] (#1,0)--(#2,0) node[right] {$x$};
  \draw[->,thick] (0,#3)--(0,#4) node[above] {$y$};
\end{tikzpicture}}
\begin{document}

% MODULE 7 PREQUIZ
\assessmenttitle{7}{Logarithms -- Prequiz}
\studentline
\directions{Four short questions. Give exact answers when reasonable and decimal approximations when requested.}
\begin{enumerate}
\item Rewrite $3^4=81$ in logarithmic form.
\answerblank{0.42in}
\item Evaluate $\log_2(32)$ and $\ln(e^3)$.
\answerblank{0.48in}
\item Expand $\log\!\left(\dfrac{x^2}{y}\right)$ using logarithm laws.
\answerblank{0.52in}
\item Solve $5^{x-1}=12$. Give an exact answer using logarithms and a decimal approximation to three places.
\answerblank{0.78in}
\end{enumerate}
\newpage

% MODULE 7 CHECKIN PAGE 1
\assessmenttitle{7}{Logarithms -- Weekly Check-In}
\studentline
\directions{Three questions. Give exact answers when reasonable and decimal approximations when requested.}
\textbf{1. Exponential and logarithmic forms.}

\begin{enumerate}[label=(\alph*),leftmargin=0.35in]
\item Rewrite $2^5=32$ in logarithmic form.
\answerblank{0.36in}
\item Evaluate $\log_3(81)$ and $\ln(e^{-2})$.
\answerblank{0.45in}
\item Expand $\log\!\left(\dfrac{x^3\sqrt{y}}{z^2}\right)$ using logarithm laws.
\answerblank{0.65in}
\end{enumerate}

\textbf{2. Solving equations with logarithms.}

\begin{enumerate}[label=(\alph*),leftmargin=0.35in]
\item Solve $4^{x-2}=11$. Give an exact answer using logarithms and a decimal approximation.
\answerblank{0.78in}
\item Solve $\ln(2x-1)=3$. Give an exact answer and check the logarithm's domain.
\answerblank{0.82in}
\end{enumerate}
\newpage

% MODULE 7 CHECKIN PAGE 2
\assessmenttitle{7}{Logarithms -- Weekly Check-In, continued}
\studentline
\textbf{3. A fading observatory signal.}

An observatory is tracking a weak radio source. The measured signal power is modeled by
\[
P(d)=1800(0.84)^d,
\]
where $P(d)$ is the signal power in arbitrary units and $d$ is the distance from the receiver in kilometers.
\begin{enumerate}[label=(\alph*),leftmargin=0.35in]
\item What is the signal power at the receiver, when $d=0$?
\answerblank{0.48in}
\item Find the signal power when the source is $5$ km away. Round to the nearest whole unit.
\answerblank{0.72in}
\item At what distance will the signal power drop to $300$ units? Set up and solve an exponential equation using logarithms. Round to three decimal places.
\answerblank{1.55in}
\end{enumerate}

\end{document}
